\section{Caps and directions in a random subspace}\label{s:subspace-moments}

The residual law reduces a block to averages of continuation functions
on a random subspace. We first bound terminal success by spherical caps.
We then compute the dependence between several directions drawn in the
same random subspace. Both block methods will use this conditional law;
their different ways of bounding its moments come later.

\subsection{Caps and terminal outputs}

\begin{lemma}[Spherical caps]\label{s:caps}
Let \(e\) be a unit vector in a \(D\)-dimensional Euclidean space.
For \(D\ge3\) and every \(\delta>0\),
\begin{equation}\label{s:sharp-cap}
 \sigma_{\R^D}\{u:\norm{u-e}\le\delta\}
 \le \frac{c_D}{2}\delta^{D-1}
 \le \frac{\sqrt D}{2}\delta^{D-1}.
\end{equation}
For \(D\ge2\) and \(0\le\theta\le\pi\),
\begin{equation}\label{s:angular-cap}
 \sigma_{\R^D}\{u:\angle(u,e)\le\theta\}
 \le(2\theta)^{D-1}.
\end{equation}
Consequently, on an affine subsphere \(z+R S(H)\) as in
\eqref{s:affine-sphere}, with radius \(R\) and linear dimension
\(D\ge2\), the angular success probability of any
fixed unit output is at most
\begin{equation}\label{s:affine-terminal}
 \min\{1,(8\epsilon/R)^{D-1}\}.
\end{equation}
The last assertion also holds after averaging any target-independent
output law.
\end{lemma}
\begin{proof}
For the chordal cap, \(t=\ip ue\) lies in an interval of length at
most \(\delta^2/2\) next to \(1\), and \(1-t^2\le\delta^2\) on that
interval. Integrating \eqref{s:coordinate-density} with \(l=1\)
gives \eqref{s:sharp-cap}. For the angular cap, the polar-angle density
is proportional to \(\sin^{D-2}\phi\) on \([0,\pi]\). Its numerator
over \([0,\theta]\) is at most \(\theta^{D-1}/(D-1)\).
Its denominator is at least \(2^{-(D-2)}\), by integration over
\([\pi/6,5\pi/6]\). This proves \eqref{s:angular-cap}.

Angular error at most \(\epsilon\) between unit vectors implies
Euclidean error at most \(\epsilon\). After projection onto \(H\)
and division by \(R\), the successful directions in
\eqref{s:affine-sphere} lie in a Euclidean ball of radius
\(\epsilon/R\). If that set is nonempty, choose one successful
direction as a center. Every other successful direction has chordal
distance at most \(2\epsilon/R\) from it. For chordal distance
\(\delta\le2\), the corresponding angle is at most
\(2\arcsin(\delta/2)\le(\pi/2)\delta\). When \(8\epsilon/R<1\),
\eqref{s:angular-cap} therefore gives
\((8\epsilon/R)^{D-1}\); when \(8\epsilon/R\ge1\), total mass is
the asserted bound. Integrating this pointwise statement over an
independent output law proves the last claim.
\end{proof}

\subsection{Directions with a common random subspace}

The joint direction law below is a spherical form of the linear
Blaschke--Petkantschin formula; see Rubin
\cite[Corollary~3.2, equation~(3.5)]{Rubin2018}.
We derive the normalized conditional density directly.

\begin{lemma}[Conditional direction law]\label{s:directions}
Let \(E\) have dimension \(N\), let \(L\) be a uniform
\(K\)-dimensional subspace, and assume integers \(1\le h<K\le N\).
Conditional on \(L\), draw \(Y_1,\ldots,Y_h\) independently with
law \(\sigma_L\). Given \(Y_1,\ldots,Y_j\), their span \(V\) has
dimension \(j\) almost surely, and \(L\) is uniform among the
\(K\)-subspaces containing \(V\). For \(1\le j<h\), the conditional
density of \(Y_{j+1}\) relative to \(\sigma_E\) is
\begin{equation}\label{s:direction-density}
 \frac{c(K,j)}{c(N,j)}
 \dist(y,V)^{-(N-K)}
 \le \dist(y,V)^{-(N-K)}.
\end{equation}
For \(j=0\), the first marginal is \(\sigma_E\). If \(K=N\), the
density in \eqref{s:direction-density} is identically one.
\end{lemma}
\begin{proof}
Fix a \(K\)-space \(L_0\subset E\). Draw independent
\(\xi_1,\ldots,\xi_h\sim\sigma_{L_0}\) and an independent uniform
orthogonal matrix \(Q\), and put \(L=QL_0\), \(Y_i=Q\xi_i\).
This produces the stated experiment. Since \(h<K\), the first \(j\)
preimages, and hence their images, are linearly independent almost
surely.

Condition on the first \(j\) preimages and images. Applying
Gram--Schmidt to the preimages gives an orthonormal basis of their
span, and the images determine its image basis. Conditional on these
images, the action of \(Q\) on the perpendicular complement is a
uniform orthogonal isometry. It follows that
\[
 L=V\oplus L',\qquad
 L'\text{ uniform of dimension }K-j\text{ in }V^\perp.
\]
This conditional law depends only on \(V\), so it persists when the
preimages are forgotten.

The next preimage is still uniform in \(L_0\). Its projection onto
the preceding span has the \(j\)-coordinate density
\eqref{s:coordinate-density} in dimension \(K\). Its complementary
direction is sent by the uniform remaining isometry to a uniform
direction in \(V^\perp\), independently of that projection. The same
decomposition of a direction with law \(\sigma_E\) uses the coordinate
density in dimension \(N\). Taking their ratio gives
\[
 \frac{c(K,j)}{c(N,j)}
 (1-\norm{P_Vy}^2)^{-(N-K)/2},
\]
which is \eqref{s:direction-density}. The integral normalizing
\(c(K,j)\) is at least that normalizing \(c(N,j)\), because
\((1-\norm t^2)^{(K-j-2)/2}\ge
 (1-\norm t^2)^{(N-j-2)/2}\) on the unit ball. Thus the prefactor
is at most one. The first marginal is uniform by rotation invariance.
When \(K=N\), \(L=E\) and the original conditional draws remain
independent uniform directions, as the formula also shows.
\end{proof}

